\section{Evaluation and Discussion}
\label{sec:evaluation}

\subsection{Performance vs. Compatibility}

The results across the three stages reveal a fundamental tension. Lightpanda delivers genuine and reproducible performance benefits in synthetic settings: approximately 2$\times$ faster request handling and 10$\times$ lower memory in controlled micro-benchmarks, and a modest but real 1.23$\times$ total pipeline speedup in a simple CI workflow. These gains are consistent across hardware architectures (Intel x86-64 and Apple M3 Max) and arise from Lightpanda's lightweight single-process design, which avoids the rendering, GPU, and network service overhead of Chrome's multi-process architecture \cite{reis2009isolating}.

However, performance advantages are irrelevant if the engine cannot execute the test suite. Stage 3 demonstrates unambiguously that Lightpanda cannot serve as a drop-in replacement for Chrome Headless in real-world E2E testing: a 75\% test failure rate on a simple e-commerce application renders the engine unusable for this purpose regardless of its speed.

\subsection{The Core Compatibility Problem}

The most critical missing capability is the absence of isolated browser contexts. The Playwright best practice of creating a fresh \texttt{BrowserContext} per test is not supported by Lightpanda, which means the standard test isolation pattern fails at the framework level before any test logic runs. This is not an edge case -- it is the default and recommended way to structure Playwright test suites \cite{playwright2024}. Until Lightpanda implements \texttt{browser.newContext()}, it is incompatible with the overwhelming majority of professionally structured Playwright suites.

The missing CDP domains and incomplete JavaScript engine coverage compound this problem. Angular applications, like many modern frameworks, rely on browser APIs (event dispatch, \texttt{MutationObserver}, layout APIs) that Lightpanda has not yet implemented. Playwright's own auto-waiting mechanism internally uses CDP calls that similarly depend on domain coverage Lightpanda does not provide.

\subsection{Critique of Lightpanda's Published Benchmark}

Lightpanda's headline performance claims -- 9$\times$ faster execution, 16$\times$ less memory -- deserve scrutiny in light of the Stage 3 results. These figures are derived from a ``crawler benchmark'' that navigates 933 pages of \textit{\url{demo-browser.lightpanda.io/amiibo/}}: a purpose-built demo page consisting of a single image and a small number of internal links. This page requires no JavaScript framework, no authentication, no dynamic DOM mutation, and no complex CSS layout -- precisely the features that expose Lightpanda's limitations in practice.

Claiming that this constitutes a benchmark across 933 ``real web pages'' is misleading. Real web applications, as demonstrated by Stage 3, rely on framework runtimes (Angular, React), client-side routing, dynamic content loading, form interactions, and authenticated sessions -- none of which the Lightpanda demo page exercises. The benchmark has no ecological validity for E2E testing workloads, and the headline figures should not be used to evaluate Lightpanda's suitability as a Chrome replacement in production pipelines.

\subsection{Where Lightpanda Is Applicable Today}

The Stage 1 and Stage 2 results suggest a narrow but legitimate use case: web crawling and simple content extraction on statically rendered or minimally interactive pages. If a test suite is limited to basic navigation and DOM visibility assertions on simple pages -- similar in complexity to the TodoMVC demo -- and does not require browser context isolation, Lightpanda can provide a real speed and memory benefit. This corresponds roughly to what its benchmark actually measures, even if the marketing language overstates the scope.

For any team running a realistic E2E test suite against a modern web application, Lightpanda in its current state (v0.3.0, beta) is not viable. The Lightpanda GitHub issue tracker openly acknowledges this, listing missing features like limited WebSocket support and numerous missing CDP features as known limitations of the current release.
